El propósito de este capítulo es establecer las políticas, responsabilidades y prácticas preventivas que permitan identificar, mitigar y controlar los efectos de los factores humanos en las actividades de mantenimiento de las aeronaves UAS de SKYFILMS.\\

Los factores humanos son aquellos elementos fisiológicos, psicológicos, organizacionales y ambientales que pueden influenciar el rendimiento del personal técnico, y por tanto, impactar la \textbf{seguridad operacional y la calidad del mantenimiento}.\\

Este procedimiento busca fomentar una cultura de seguridad y desempeño humano, promoviendo la conciencia situacional, la comunicación efectiva y el autocontrol del personal en todas las etapas del proceso técnico.

\section{Alcance}

Este capítulo aplica a todo el personal involucrado en las actividades de mantenimiento, inspección, control técnico, almacenamiento y gestión documental dentro del área técnica de SKYFILMS.\\

El procedimiento se aplica tanto a operaciones internas como externas, en locaciones de rodaje o zonas de operación en campo, donde las condiciones de trabajo puedan afectar el desempeño humano.

\section{Concepto de Factores Humanos}
Los factores humanos en el mantenimiento se refieren a las interacciones entre las personas, las herramientas, el entorno y los procedimientos. Su correcta gestión busca reducir la probabilidad de error y mejorar la eficiencia del equipo.\\

Entre los principales factores que pueden influir negativamente en el desempeño del personal se incluyen:

\begin{itemize}
    \item Fatiga física o mental.

    \item Presiones operativas o de tiempo.

    \item Falta de comunicación efectiva.

    \item Ambientes de trabajo inadecuados (iluminación, ruido, temperatura).

    \item Distracciones o interrupciones durante tareas críticas.

    \item Deficiencias de entrenamiento o experiencia.

    \item Actitudes complacientes o exceso de confianza.
\end{itemize}

Reconocer y gestionar estos factores es esencial para mantener altos niveles de seguridad y calidad en las tareas de mantenimiento.

\section{Principios de Gestión de Factores Humanos}
SKYFILMS adopta una política activa para el control de los factores humanos, basada en los siguientes principios:

\begin{enumerate}
    \item \textbf{Seguridad ante todo:} Ninguna tarea debe ejecutarse bajo condiciones de cansancio extremo, estrés o presión operativa que comprometa la precisión técnica.

    \item \textbf{Conciencia situacional:} El personal debe mantener una percepción constante del entorno, los riesgos y las consecuencias de cada acción realizada.

    \item \textbf{Trabajo en equipo y comunicación efectiva:} Las instrucciones deben transmitirse con claridad y verificarse mediante confirmaciones o registros escritos.

    \item \textbf{Cumplimiento estricto de procedimientos:} Toda tarea debe realizarse conforme a los manuales técnicos y procedimientos del MCM, MO y SMS sin omisiones o improvisaciones.

    \item \textbf{Gestión del error humano:} Los errores deben ser reportados de forma abierta y sin sanción punitiva, para analizar sus causas y prevenir su repetición.
\end{enumerate}

Estos principios son transversales a todas las actividades y reflejan el compromiso de SKYFILMS con la seguridad operacional y la integridad técnica.

\section{Identificación de Factores Humanos}
El Técnico de UAS es responsable de supervisar que los factores humanos sean evaluados y gestionados de manera proactiva durante todas las fases del mantenimiento.

Las áreas clave donde se deben identificar riesgos asociados son:

\begin{enumerate}
    \item \textbf{Planificación y asignación de tareas:}
        \begin{itemize}
            \item Verificar que la carga de trabajo esté equilibrada y que se cuente con tiempo suficiente para ejecutar sus tareas.

            \item Asegurar que las tareas críticas sean revisadas por un segundo.
        \end{itemize}

    \item \textbf{Condiciones del entorno de trabajo:}
        \begin{itemize}
            \item Mantener las áreas limpias, iluminadas, ventiladas y con temperatura controlada.

            \item Evitar distracciones o interrupciones durante labores de precisión.
        \end{itemize}

    \item \textbf{Capacitación y competencia:}
        \begin{itemize}
            \item Garantizar que todo el personal técnico esté adecuadamente entrenado en el modelo de UAS sobre el cual trabaja.

            \item Promover actualizaciones periódicas en temas de factores humanos, comunicación y seguridad operacional.
        \end{itemize}

    \item \textbf{Estado físico y mental del personal:}
        \begin{itemize}
            \item Evaluar la fatiga, estrés o condiciones personales que puedan afectar la concentración y el juicio técnico.

            \item Permitir pausas razonables durante turnos prolongados o tareas de alta exigencia.
        \end{itemize}

    \item \textbf{Supervisión y liderazgo:}
        \begin{itemize}
            \item Fomentar un entorno de respeto y colaboración.

            \item El líder técnico debe observar y corregir comportamientos inseguros, guiando al personal hacia la mejora continua.
        \end{itemize}

\end{enumerate}

\section{Estrategias de Mitigación de Factores Humanos}
SKYFILMS implementa diversas estrategias preventivas para reducir la influencia negativa de los factores humanos en sus operaciones técnicas:

\begin{itemize}
    \item \textbf{Gestión de fatiga:} asignación de horarios equilibrados y descanso obligatorio entre turnos.

    \item \textbf{Capacitación continua:} programas periódicos sobre seguridad, gestión del error humano y toma de decisiones.

    \item \textbf{Revisión cruzada de tareas críticas:} toda inspección o mantenimiento mayor debe ser verificada por un segundo.

    \item \textbf{Reportes de seguridad operacional:} canal confidencial de notificación voluntaria de errores o eventos.

    \item \textbf{Evaluación del entorno de trabajo:} revisiones mensuales de condiciones de iluminación, orden, ruido y ergonomía.

    \item \textbf{Retroalimentación técnica:} reuniones periódicas para analizar incidentes, compartir aprendizajes y fortalecer la cultura de seguridad.
\end{itemize}

Estas medidas contribuyen a mantener la confiabilidad y la seguridad operacional en todas las fases del servicio.

\section{Interacción con la Máquina, el Entorno y las Personas}
\begin{itemize}
    \item \textbf{Persona y máquina:} el Técnico UAS trabaja con los manuales del fabricante a la vista y sigue las tareas en orden, con una lista de chequeo, para no depender de la memoria. Los mensajes de error de la aeronave o del control se anotan tal como aparecen.
    \item \textbf{Ergonomía y entorno:} se trabaja sentado, con la aeronave a la altura de las manos, con luz LED suficiente, sin ruido ni interrupciones y con el tapete antiestático en su sitio. Las piezas pequeñas se guardan en recipientes para que no se pierdan.
    \item \textbf{Fiabilidad de la inspección visual:} una inspección visual detallada no se hace durante más de 30 minutos seguidos; luego se descansa unos minutos. Las tareas críticas (motores, baterías, estructura) las revisa una segunda persona, que puede ser el Jefe de Pilotos.
    \item \textbf{Comunicación entre personas:} al entregar la aeronave, el Técnico UAS y el piloto repasan por escrito qué se hizo y qué queda pendiente, en el F-005 o el F-022. Nada se entrega solo de palabra.
    \item \textbf{Diseño de los documentos:} los formatos son cortos, usan casillas y lenguaje sencillo y tienen guía de diligenciamiento. Si una persona se equivoca de forma repetida en un formato, se corrige el formato y no solo a la persona.
\end{itemize}

\section{Función de la Administración y Vigilancia Reglamentaria}
\begin{itemize}
    \item El Gerente General da el tiempo, las herramientas y la capacitación para trabajar sin prisa, y no presiona para liberar una aeronave sin terminar el mantenimiento.
    \item El Gerente SMS revisa cada trimestre los reportes de mantenimiento, los errores repetidos y la carga de trabajo, y propone acciones.
    \item Cada año se verifica que el trabajo se hace como dicen el RAC 100 y este Manual, y los hallazgos se registran y se cierran.
    \item Los errores se reportan sin castigo cuando se informan a tiempo; la negligencia o el ocultamiento sí se sancionan.
\end{itemize}

\section{Responsabilidades}

Las responsabilidades en la gestión de los factores humanos dentro de SKYFILMS se distribuyen de la siguiente manera:

\begin{itemize}
    \item \textbf{Gerente General:} Establece las políticas corporativas relacionadas con seguridad humana y proporciona los recursos necesarios para su implementación.

    \item \textbf{Técnico de UAS:}
        \begin{itemize}
            \item Supervisa el cumplimiento del programa de control de factores humanos, asegura la capacitación periódica del personal y promueve un entorno laboral seguro.

            \item Aplica los procedimientos, realiza las evaluaciones periódicas, mantiene los registros y reporta cualquier riesgo humano identificado.

            \item Deben trabajar de forma segura, cumplir las pausas establecidas, reportar condiciones de fatiga o estrés y comunicar cualquier incidente que comprometa la seguridad.
        \end{itemize}

\end{itemize}

Cada rol es esencial para prevenir errores humanos y garantizar que las tareas de mantenimiento se ejecuten con la máxima precisión y responsabilidad.

\section{Disposiciones Finales}
La gestión de los factores humanos constituye una parte esencial del Sistema de Gestión de Seguridad Operacional (SMS) de SKYFILMS. La comprensión y control de estos factores refuerzan la prevención de incidentes y errores técnicos, promoviendo la eficiencia, el profesionalismo y la seguridad en cada proceso.\\

El incumplimiento de las políticas aquí descritas será considerado una no conformidad técnica y administrativa, sujeta a acciones correctivas según los procedimientos internos de la organización.\\

Mediante la aplicación de este capítulo, SKYFILMS reafirma su compromiso con la seguridad operacional, el bienestar del personal técnico y la calidad total en el mantenimiento de sus aeronaves UAS.